\scanpage{226}
\margin{cet \unsure{énoncé} a un sens et est valable \unsure{en supposant} $B$ karoubienne \unsure{ici}}
\[
F(x)\simeq\operatorname{Hom}^*_k(\alpha^*(x),
\underset{\substack{\Vert\,\mathrm{déf}\\F(\mathbb C^*)}}{C_F^*})
\quad\deleted{\unsure{\in\widehat B}}
\]
La définition des $\operatorname{Hom}^*$ se fait comme dans la prop. 3.

\deleted{
\diagram{\begin{tikzcd}[ampersand replacement=\&]
F(x)\otimes\alpha^{[n]}(x)\arrow[r,Rightarrow] \& C_F^{[n]}\arrow[d,equal]\\
\operatorname{Hom}(x,\mathbb C^{[n]})\arrow[u,equal] \& F(\mathbb C^{[n]})
\end{tikzcd}}
\margin{par définition des $\alpha^*$ et \deleted{$\unsure{DP}$} \unsure{adjoint|adéquat}
\deleted{$\unsure{C_*}$} à \unsure{l'application de} $DP$ \unsure{aux} \unsure{formules}. \unsure{c.d.}}
\unsure{procéder} comme pour la prop. 3 pour définir la flèche dans $\widehat B$,
et pour montrer que c'est un iso, on se ramène au cas $x$ de la forme
$\mathbb C^i_{\unsure{\bullet}}$, il faut prouver que
\[
F(\mathbb C^i)=C^i\longrightarrow
\operatorname{Hom}^*_k(\deleted{\unsure{A}}\mathbb C^{\unsure{*|i}}_{\unsure{\bullet}})
\]
}
En effet, on trouve tout de suite que si
$K^\bullet\in\mathcal M_k^\bullet$, $C^\bullet\in B^\bullet$,
$K^*=DP(K^\bullet)$ et $C^*=DP(\unsure{K^\bullet|C^\bullet})$,
alors on a un iso. can. \unsure{bifonctoriel} dans $\widehat B$
\[
DP:\operatorname{Hom}_k^\bullet(K^\bullet,C^\bullet)
\simeq\operatorname{Hom}_k^*(K^*,C^*),
\]
et on conclut par prop 3.

\deleted{
\uline{Prop 6} $^{\unsure{0|o}}$ Le foncteur
\[
ND:\mathrm{Ab}_{k,*}\longrightarrow\mathcal M_{k,\bullet}
\]
est donné par
\[
(ND)(L_*)\simeq\operatorname{Hom}_{k,*}(C_*^\bullet,L_*)
\]
i.e.
\[
\underline{ND}(L_*)_i\simeq\operatorname{Hom}_{k,*}
(\underset{\substack{\Vert\\\Lambda^{i+1}\Phi_*}}{C_*^i},L_*)
\quad\text{pour }i\geq0
\]
l'hom \inserted{bord}
$ND(L_*)_i\mathbin{-}ND(L_*)_{i-1}$ $(i\geq1)$ est le transposé de
$C_*^{i-1}\longrightarrow C_*^i$.
}
\margin{225}
